\begin{frame}\frametitle{Objetivos del tema}
  \alert{Establecer la terminología adecuada y repasar conceptos:}
  \vspace{0.1cm}
  \begin{itemize}
    \item \alert{Señales}
    \begin{itemize}
      \item Tipos
      \item Operaciones
      \item Ejemplos
    \end{itemize}
    \item \alert{Sistemas}
    \begin{itemize}
      \item Tipos
      \item Propiedades de los sistemas lineales temporalmente invariantes
    \end{itemize}
  \end{itemize}
\end{frame}


\section{Definiciones y notación}

\begin{frame}
  \frametitle{Definiciones}

  \begin{figure}
    \centering
    \small
    \begin{tikzpicture}
      \node[bloque] (T) {Sistema};
      \node (x) [left=of T] {Señal};
      \node (y) [right=of T] {Señal};
      \draw[arrow,->] (x.east) -- (T.west);
      \draw[arrow,->] (T.east) -- (y.west);
    \end{tikzpicture}
  \end{figure}
  \vspace{-0.3cm}
  \begin{itemize}
    \item \alert{Señal:}
    \begin{itemize}
      \item Función que transporta algún tipo de información sobre la variabilidad/evolución de una magnitud física respecto a una/varias variables independientes.
      \vspace{0.1cm}
      \item Ejemplos: Señales de audio, imagen, video, tensión, presión, ...
    \end{itemize}
    \item \alert{Sistema}
    \begin{itemize}
      \item Definición:
      \begin{itemize}
        \item Entidad que se excita mediante una señal de entrada y produce una señal de salida.
        \item Representa la transformación que un sistema físico realiza sobre una señal.
      \end{itemize}
      \item Ejemplos: filtro paso bajo, conversor A/D, ...
    \end{itemize}
  \end{itemize}
\end{frame}

\begin{frame}
  \frametitle{Notación}
  \begin{figure}
    \centering
    \small
    \begin{tikzpicture}
      \node[bloque,minimum width=15mm] (T) {T};
      \node (x) [left=of T] {x(t)};
      \node (y) [right=of T] {y(t)=T[x(t)]};
      \draw[arrow,->] (x.east) -- (T.west);
      \draw[arrow,->] (T.east) -- (y.west);
    \end{tikzpicture}
  \end{figure}

  
  \begin{block}{}
  \begin{itemize}
    \item Señales: $x,y,u,...$ (letras minúsculas)
    \begin{itemize}
      \item variable independiente: $t$, $n$
      \item $y(t)$ - tiempo continuo
      \item $y(n), y[n]$ - secuencia
    \end{itemize}
    \item Sistemas: $T,P,S,...$ (letras mayúsculas)
    \begin{itemize}
      \item $y(t)=T[x(t)]$
      \item $y[n]=T[x[n]]$
    \end{itemize}
  \end{itemize}
  \end{block}
\end{frame}
